%% Dans ce fichier, les �ventuels accents sont cod�s
%% en ISO-8859-15 (aka latin9)
%% La ligne suivante donne trois exemples 
%% e accent aigu : �, e accent grave : �, e-dans-l'o : �
%%%%%ENTREES%%%%
\begin{tabarg}
\multicolumn{3}{l}{\titretab {Entr\'ees}}\\
\hline
$X$ & Vecteur, \texttt{float} & Op�rande (unit� indiff�rente)\\
$C_0$ & Coeff., \texttt{float} & Coefficient multiplicateur du logarithme\\
$C_1$ & Coeff., \texttt{float} & Coefficient \\
$C_2$ & Coeff., \texttt{float} & Coefficient \\
$C_3$ & Coeff., \texttt{float} & Coefficient constant\\
$Unit\_res$ & Coeff., \texttt{String} & Unit� du r�sultat\\
\$Arg7 & Coeff., \texttt{String} & description mono-ligne du r�sultat (en fran�ais)\\
\$Arg8 & Coeff., \texttt{String} & description mono-ligne du r�sultat (en anglais)\\
\end{tabarg}

\medskip
%%%%%SORTIES%%%%
\begin{tabarg}
\multicolumn{3}{l}{\titretab {Sorties}}\\
\hline
$Z$ & Vecteur, \texttt{float} & R�sultat (unit� indiqu�e par $Unit\_res$)\\
\end{tabarg}
